\section{Constraint Modeling}
\label{sec:stage3}
% ============================================================

\subsection{Structured Constraint Representation}
\textit{[Para: All constraints are stored in a nested structured format (JSON):
(a) univariate distributions (type, parameters, target column);
(b) decision-tree logical constraints (if--then rules over column values, including
cross-table aggregations such as SUM, MAX);
(c) multivariate distribution groups (\S\ref{subsec:multivariate}).
This representation enables downstream deterministic validation and code generation.]}

\subsection{Constraint Extraction and Feasibility Verification}
\textit{[Para: The global schema graph is sharded into connected components.
A constraint-modeling agent extracts structured constraints per component in parallel.
A mathematics agent then verifies distribution feasibility (e.g., checking that parameter
values are in the valid domain for the specified distribution family).
Decision-tree feasibility is checked deterministically (Algorithm~\ref{alg:dt-feasibility}).]}

% Algorithm 7: Deterministic Decision-Tree Feasibility Check (sub-step)
\begin{algorithm}[t]
\caption{Deterministic Decision-Tree Feasibility Check}
\label{alg:dt-feasibility}
\begin{algorithmic}[1]
  \Require Set of decision-tree constraints $T \subseteq \mathcal{C}$
  \Ensure Feasibility verdict per constraint; infeasibility report if any
  \State \textit{[Placeholder: per constraint---enumerate leaf conditions,
         check contradictory conjunctions (e.g., $x > 750$ AND $x < 300$),
         verify outcome type compatibility, check referenced columns exist in $\mathcal{S}$;
         flag infeasible constraints for LLM repair]}
\end{algorithmic}
\end{algorithm}

\subsection{Multivariate Distribution Support}
\label{subsec:multivariate}
\textit{[Para --- PLACEHOLDER (not fully fleshed out):
Extension beyond univariate constraints to capture cross-column correlations.
Planned approach: [TBD --- likely copula-based joint sampling (Gaussian or vine copulas)
or a conditional chain representation where later columns in the DAG topological order
are sampled conditionally on earlier ones].
Describe: (i) the structured representation for multivariate groups in $\mathcal{C}$;
(ii) how multivariate group nodes are inserted into the dependency graph $\mathcal{G}$;
(iii) how the code generator handles a multivariate group node.
This section will expand substantially once the approach is finalised.]}

\subsection{Dependency Graph Construction}
\textit{[Para: From FK links and structured constraints, a column-level DAG $\mathcal{G}$
is built in \emph{reverse-causal} order: constrained outcome columns
(e.g., \texttt{interest\_rate}) are positioned as roots; their independent drivers
(e.g., \texttt{credit\_score}, \texttt{loan\_amount}) are leaves.
Multivariate group nodes are inserted as composite nodes (TBD).
Cycle detection is performed via topological sort; if a cycle is detected, a patch agent
updates $\mathcal{C}$ to break the cycle and the graph is rebuilt.]}

% Algorithm 8: Deterministic DAG Construction (sub-step)
\begin{algorithm}[t]
\caption{Deterministic Dependency DAG Construction}
\label{alg:dag}
\begin{algorithmic}[1]
  \Require Schema $\mathcal{S}$, structured constraints $\mathcal{C}$
  \Ensure Column-level DAG $\mathcal{G}$
  \State \textit{[Placeholder: one node per column; FK edges reverse-causal;
         constraint edges outcome $\leftarrow$ driver; multivariate group nodes (TBD);
         topological sort; cycle $\to$ patch agent updates $\mathcal{C}$, rebuild]}
\end{algorithmic}
\end{algorithm}

% Figure 5: Dependency DAG
\begin{figure}[t]
  \centering
  \vspace{1.5cm}
  \textit{[Figure 5 placeholder: Column-level DAG for Credit Risk.
  Solid nodes: \texttt{credit\_score} (leaf), \texttt{loan\_amount} (leaf),
  \texttt{interest\_rate} (root).
  Solid directed edges: \texttt{credit\_score} $\to$ \texttt{interest\_rate},
  \texttt{loan\_amount} $\to$ \texttt{interest\_rate}.
  FK edges shown as dashed.
  Dashed composite node group for a multivariate dependency (placeholder).
  Node colouring: deterministic-path nodes (blue), LLM-path nodes (orange).]}
  \vspace{1.5cm}
  \caption{Dependency DAG for the Credit Risk running example.
           \texttt{interest\_rate} is the reverse-causal root; its drivers are leaves.
           Node colour indicates the code-generation path taken in Stage 4.}
  \label{fig:dag}
\end{figure}

% Algorithm 9: Stage 3 summary (calls Algs 7 & 8)
\begin{algorithm}[t]
\caption{Stage 3: Constraint Modeling}
\label{alg:stage3}
\begin{algorithmic}[1]
  \Require Schema $\mathcal{S}$, fact set $\mathcal{F}$
  \Ensure Structured constraints $\mathcal{C}$, dependency DAG $\mathcal{G}$
  \State \textit{[Placeholder: shard $\mathcal{S}$ into connected components
         $\to$ parallel constraint-modeling agents $\to$ math agent distribution feasibility
         $\to$ DTFeasibilityCheck (Alg.~\ref{alg:dt-feasibility})
         $\to$ BuildDAG (Alg.~\ref{alg:dag})
         $\to$ return $\mathcal{C}, \mathcal{G}$]}
\end{algorithmic}
\end{algorithm}

